% =============================================================================
\section{Análisis de Fourier}
% =============================================================================

\begin{frame}{Serie de Fourier}
  Toda señal periódica $x(t)$ de período $T_0$ que cumple las condiciones de Dirichlet
  admite el desarrollo
  \begin{align*}
    x(t) &= \sum_{k=-\infty}^{\infty} c_k\, \e^{\j k \omega_0 t},
    & \omega_0 &= \frac{2\pi}{T_0} \\
    % En align* usar \onslide por celda (\pause no funciona bien: amsmath mide dos veces)
    \onslide<2->{c_k} &\onslide<2->{= \frac{1}{T_0} \int_{\langle T_0 \rangle} x(t)\, \e^{-\j k \omega_0 t} \dd t}
  \end{align*}
  \begin{teorema}<3->[Parseval para señales periódicas]
    \[
      P_x = \frac{1}{T_0} \int_{\langle T_0 \rangle} \abs{x(t)}^2 \dd t
          = \sum_{k=-\infty}^{\infty} \abs{c_k}^2
    \]
  \end{teorema}
\end{frame}

\begin{frame}{Transformada de Fourier}
  \begin{definicion}
    \vspace{-0.5\baselineskip}
    \begin{align*}
      X(\j\omega) = \TF{x(t)} &= \int_{-\infty}^{\infty} x(t)\, \e^{-\j\omega t} \dd t \\
      x(t) = \iTF{X(\j\omega)} &= \frac{1}{2\pi} \int_{-\infty}^{\infty} X(\j\omega)\, \e^{\j\omega t} \dd \omega
    \end{align*}
  \end{definicion}
  Se escribe $x(t) \fourier X(\j\omega)$.
  \begin{columns}[T]
    \begin{column}{0.48\textwidth}
      \begin{itemize}
        \item $X(\j\omega) = \abs{X(\j\omega)}\, \e^{\j\angle X(\j\omega)}$
        \item Espectro de magnitud y de fase
      \end{itemize}
    \end{column}
    \begin{column}{0.48\textwidth}
      \begin{itemize}
        \item $x(t)$ real $\Rightarrow X(-\j\omega) = \conj{X}(\j\omega)$
        \item Simetría hermítica
      \end{itemize}
    \end{column}
  \end{columns}
\end{frame}

\begin{frame}{Pares y propiedades}
  \begin{columns}[T]
    \begin{column}{0.47\textwidth}
      \centering
      \renewcommand{\arraystretch}{1.25}
      \begin{tabular}{@{}cc@{}}
        \toprule
        $x(t)$ & $X(\j\omega)$ \\
        \midrule
        $\impulso{t}$ & $1$ \\
        $\e^{-at}\escalon{t}$ & $\dfrac{1}{a + \j\omega}$ \\
        $\rect\!\left(\dfrac{t}{\tau}\right)$ & $\tau \sinc\!\left(\dfrac{\omega\tau}{2\pi}\right)$ \\
        $\e^{\j\omega_0 t}$ & $2\pi\, \impulso{\omega - \omega_0}$ \\
        $\cos(\omega_0 t)$ & $\pi\left[\impulso{\omega - \omega_0} + \impulso{\omega + \omega_0}\right]$ \\
        \bottomrule
      \end{tabular}
    \end{column}
    \begin{column}{0.5\textwidth}
      \centering\small
      \renewcommand{\arraystretch}{1.25}
      \begin{tabular}{@{}lc@{}}
        \toprule
        Propiedad & Resultado \\
        \midrule
        Desplazamiento & $x(t - t_0) \fourier \e^{-\j\omega t_0} X(\j\omega)$ \\
        Modulación & $\e^{\j\omega_0 t} x(t) \fourier X(\j(\omega - \omega_0))$ \\
        Derivación & $\dfrac{\dd x}{\dd t} \fourier \j\omega\, X(\j\omega)$ \\
        Convolución & $x \conv h \fourier X(\j\omega)\, H(\j\omega)$ \\
        Producto & $x\,y \fourier \dfrac{1}{2\pi}\, X \conv Y$ \\
        \bottomrule
      \end{tabular}
    \end{column}
  \end{columns}
  \vfill
  {\footnotesize Con $\sinc(x) = \dfrac{\sin(\pi x)}{\pi x}$. Tablas completas en \cite{oppenheim1997}.}
\end{frame}

\begin{frame}{Espectro de un pulso rectangular}
  \begin{columns}[c]
    \begin{column}{0.36\textwidth}
      \[
        \rect\!\left(\frac{t}{\tau}\right) \fourier \frac{2 \sin(\omega\tau/2)}{\omega}
      \]
      \begin{itemize}
        \item Ceros en $\omega = \dfrac{2\pi k}{\tau}$, $k \neq 0$
        \item<2-> Pulso más angosto $\Rightarrow$ \alert{espectro más ancho}
      \end{itemize}
    \end{column}
    \begin{column}{0.6\textwidth}
      \centering
      \begin{tikzpicture}
        \begin{axis}[espectro, width=\linewidth, height=5.2cm,
            ylabel={$X(\j\omega)$}, xmin=-20, xmax=20, ymin=-0.35, ymax=1.1,
            xtick={-12.566,-6.283,6.283,12.566},
            xticklabels={$-\frac{4\pi}{\tau}$,$-\frac{2\pi}{\tau}$,$\frac{2\pi}{\tau}$,$\frac{4\pi}{\tau}$},
            ytick={1}, yticklabels={$\tau$},
            legend style={font=\scriptsize, draw=none, at={(1,1)}, anchor=north east}]
          \addplot[senal continua, samples=400, domain=-20:20] {2*sin(deg(x/2))/x};
          \addlegendentry{$\tau = 1$}
          \only<2->{\addplot[energia, very thick, samples=400, domain=-20:20] {2*sin(deg(x*0.25))/x};
                    \addlegendentry{$\tau = 0{,}5$}}
        \end{axis}
      \end{tikzpicture}
    \end{column}
  \end{columns}
\end{frame}
